%% 36장 수치 적분 II
\chapter{수치 적분 II}
\label{ch:quadrature-2}

\abstract{앞 장의 구적 규칙은 매끄러운 함수를 정해진 구간에서 적분한다. 실제 적분에는 끝점의 특이점, 좁은 봉우리, 끝없는 구간, 빠른 진동이 섞여 있다. 이 장에서는 오차가 큰 구간을 골라 나누는 전역 적응 구적과, 나누면서 얻은 근사들의 수열을 외삽해 특이점의 효과를 지우는 Wynn의 $\varepsilon$-알고리즘을 다룬다. Aitken의 $\Delta^2$ 방법이 기하급수적 오차를 정확히 지운다는 것을 증명하고 그것을 일반화한 Shanks 변환과 그 점화식을 본다. 진동하는 함수와 대수적 특이점에는 무게 함수를 따로 다루는 규칙을, 끝점 특이점에는 변수를 바꾸어 피적분함수를 이중 지수적으로 줄이는 tanh-sinh 구적을 쓴다. 마지막으로 여러 겹의 적분과 무한 급수의 합을 다루고 stardust가 scipy와 같은 결과를 내려고 tanh-sinh의 노드 표를 어떤 함수로 계산해야 했는지 본다.}

\section{적응 구적}
\label{sec:adaptive}

\subsection{전역 적응 분할}

함수가 어떤 곳에서는 순하고 어떤 곳에서는 사나우면, 구간 전체에 같은 밀도의 노드를 두는 것은 낭비이다. \term{적응 구적}{adaptive quadrature}은 구간마다 앞 장의 Gauss--Kronrod 쌍으로 값과 오차의 어림을 구하고 오차가 큰 구간만 나눈다. 구간을 하나씩 재귀적으로 나누며 구간마다 할당된 허용오차를 맞추는 국소 방식과, 모든 구간을 힙에 넣어 두고 오차가 가장 큰 구간을 나누는 전역 방식이 있다. QUADPACK은 전역 방식을 쓴다(\cref{alg:global-adaptive}).

\begin{algorithm}[t]
\begin{algorithmic}[1]
\Require 피적분함수 $f$, 구간 $[a,b]$, 허용오차 $\mathrm{atol}$, $\mathrm{rtol}$, 구간 수 상한 $L$
\Ensure 적분값 $I$와 오차의 어림 $E$
\State $[a,b]$에 Kronrod 규칙을 써서 $(I_1,E_1)$을 얻고 힙에 넣는다
\While{$E=\sum_jE_j>\max(\mathrm{atol},\mathrm{rtol}\cdot|I|)$이고 구간 수 $<L$}
  \State 힙에서 오차가 가장 큰 구간 $[c,d]$를 꺼낸다
  \State 두 반쪽에 각각 Kronrod 규칙을 써서 $(I',E')$, $(I'',E'')$를 얻는다 \Comment{$f$를 $2\times21$번}
  \State $I\gets I-I_{[c,d]}+I'+I''$, 두 반쪽을 힙에 넣는다
\EndWhile
\State \Return $I$, $E$
\end{algorithmic}
\caption{전역 적응 구적}
\label{alg:global-adaptive}
\end{algorithm}

전역 방식은 허용오차를 구간마다 나누어 줄 필요가 없고 어떤 구간이 전체 오차를 지배하는지를 언제나 안다는 장점이 있다. Gander와 Gautschi는 적응 구적의 멈춤 판정이 반올림 오차 수준에서 어떻게 실패하는지를 자세히 분석했다\cite{gander2000}. 앞 장에서 본 대로 어떤 적응 방법도 처음 노드들이 놓친 봉우리를 찾을 수는 없다.

\subsection{특이점과 외삽}

끝점에 대수적 특이점이 있으면, 예컨대 $[0,1]$의 $x^{-1/2}$이면, 적응 분할은 특이점을 품은 맨 왼쪽 구간을 끝없이 반으로 나누게 된다. 그 구간의 오차는 길이를 반으로 줄일 때마다 일정한 비율로만 준다. 그러나 이 규칙성은 기회이기도 하다. 구간을 나눌 때마다 얻는 전체 적분의 근사 $I_1,I_2,\dots$는 오차가 대략 기하급수적으로 주는 수열이므로, 다음 절의 수열 가속으로 극한을 외삽할 수 있다. QUADPACK의 QAGS는 이렇게 적응 분할에 Wynn의 $\varepsilon$-알고리즘을 붙인 것이다\cite{dedoncker1978,piessens1983quadpack}. scipy의 \code{quad}는 $[0,1]$의 $x^{-1/2}$을 구간 $6$개, 함수 호출 $231$번으로 오차 $6.7\times10^{-16}$까지 적분한다.

\section{수열 가속}
\label{sec:acceleration}

\subsection{Aitken의 방법}

수열 $S_n\to S$의 오차가 정확히 기하급수적, 곧 $S_n=S+c\lambda^n$이면 세 항으로 $S$를 알 수 있다. 이웃한 차분 $\Delta S_n=S_{n+1}-S_n=c\lambda^n(\lambda-1)$의 비가 $\lambda$이기 때문이다. 이것을 정리한 것이 \term{Aitken의 $\Delta^2$ 방법}{Aitken's delta-squared process}이다\cite{aitken1926}.

\begin{proposition}
\label{prop:aitken}
\[
A_n=S_{n+2}-\frac{(\Delta S_{n+1})^2}{\Delta S_{n+1}-\Delta S_n}
\]
로 두면, $S_n=S+c\lambda^n$ $(\lambda\ne1,\ c\ne0)$인 수열에서 $A_n=S$이다. 또 $S_n-S=(c+\epsilon_n)\lambda^n$, $c\ne0$, $\epsilon_n\to0$, $0<|\lambda|<1$이면 $(A_n-S)/(S_{n+2}-S)\to0$이다.
\end{proposition}

\begin{proof}
첫 주장은 대입하면 된다. $\Delta S_{n+1}=c\lambda^{n+1}(\lambda-1)$, $\Delta S_{n+1}-\Delta S_n=c\lambda^n(\lambda-1)^2$이므로 빼는 항은 $c\lambda^{n+2}$이고 $A_n=S$이다. 둘째 주장에서는 $e_n=S_n-S$로 두고 $A_n-S=\bigl(e_{n+2}e_n-e_{n+1}^2\bigr)/(e_{n+2}-2e_{n+1}+e_n)$으로 쓴다. 분자는 $\lambda^{2n+2}\bigl((c+\epsilon_{n+2})(c+\epsilon_n)-(c+\epsilon_{n+1})^2\bigr)=\lambda^{2n+2}o(1)$이고, 분모는 $\lambda^n\bigl(c(\lambda-1)^2+o(1)\bigr)$이다. 그래서 $A_n-S=o(\lambda^{n+2})$이다.
\end{proof}

Aitken의 방법은 변환한 수열에 다시 쓸 수 있다. $\log2=\sum_{k\ge1}(-1)^{k+1}/k$의 부분합은 $20$항에서도 오차가 $2.4\times10^{-2}$이지만 그 $20$개의 부분합에 Aitken 변환을 세 번 거듭하면 오차가 $1.7\times10^{-5}$, $4.3\times10^{-8}$, $2.2\times10^{-10}$으로 준다. 교대급수의 오차는 $\lambda=-1$에 가까운 기하급수적 꼴이기 때문이다.

\subsection{Shanks 변환과 $\varepsilon$-알고리즘}

오차가 기하급수 $k$개의 합 $S_n=S+\sum_{j=1}^kc_j\lambda_j^n$이면 $2k+1$개의 항으로 $S$를 정할 수 있다. 미지수 $S,c_j,\lambda_j$를 푸는 것은 비선형이지만 결과는 Hankel 행렬식의 비로 닫힌 꼴로 쓸 수 있고 이것을 \term{Shanks 변환}{Shanks transformation} $e_k(S_n)$이라 한다\cite{shanks1955}. $k=1$이 Aitken의 방법이다. Wynn은 행렬식을 계산하지 않고 이 변환을 얻는 점화식을 발견했다\cite{wynn1956}.
\[
\varepsilon_{-1}^{(n)}=0,\qquad\varepsilon_0^{(n)}=S_n,\qquad\varepsilon_{k+1}^{(n)}=\varepsilon_{k-1}^{(n+1)}+\frac1{\varepsilon_k^{(n+1)}-\varepsilon_k^{(n)}}
\]
짝수 열 $\varepsilon_{2k}^{(n)}$이 $e_k(S_n)$이고 홀수 열은 계산을 위한 보조량이다. 이것이 \term{$\varepsilon$-알고리즘}{epsilon algorithm}이다. 30장의 Padé 근사와 깊은 관계가 있어서 거듭제곱 급수의 부분합에 쓰면 $\varepsilon_{2k}^{(n)}$이 Padé 근사 $[n+k/k]$가 된다. 짝수 열 하나는 분모의 차수가 $k$로 고정된 Padé 표의 한 줄이고, $\varepsilon$-표의 대각선이 Padé 표의 대각선이다. 위의 $\log2$의 부분합 $12$개로 표를 채우면 열 $2,4,6,8,10$의 마지막 값의 오차가 $8.1\times10^{-5}$, $9.5\times10^{-7}$, $3.7\times10^{-8}$, $3.9\times10^{-9}$, $1.0\times10^{-9}$이다.

수열 가속은 만능이 아니다. 오차의 구조를 가정하고 그 가정을 외삽하는 것이어서 가정이 맞지 않는 수열에서는 아무 소용이 없거나 오히려 해롭다. 차분의 나눗셈은 소거를 일으키므로, 수열이 이미 반올림 수준으로 수렴한 뒤에는 표가 잡음을 증폭한다. QUADPACK은 이런 경우를 판정해 외삽을 멈추는 장치를 둔다. 수열 가속의 이론은 Sidi\cite{sidi2003}, Brezinski와 Redivo Zaglia\cite{brezinski1991}의 책이 자세하다.

\section{무게 함수와 무한 구간}
\label{sec:weights-infinite}

\subsection{진동과 대수적 특이점}

피적분함수가 $f(x)\cos\omega x$처럼 빠르게 진동하는 인자와 매끄러운 인자의 곱이면, 진동 인자를 무게 함수로 떼어 내고 매끄러운 $f$만 다항식으로 근사하는 편이 낫다. QUADPACK의 QAWO는 $f$를 앞 장의 Clenshaw--Curtis처럼 Chebyshev 다항식으로 근사하고 Chebyshev 다항식과 $\cos\omega x$의 곱의 적분, 곧 \term{변형 모멘트}{modified moments}를 점화식으로 정확히 계산한다\cite{piessens1983quadpack}. 같은 생각으로 QAWS는 끝점의 $(x-a)^\alpha(b-x)^\beta$와 로그 인자를, QAWC는 Cauchy 주값의 $1/(x-c)$를 무게로 다룬다. scipy의 \code{quad}에서 무게 $x^{-1/2}$을 지정해 $[0,1]$의 $1$을 적분하면 함수를 $40$번 불러 정확한 값 $2$를 낸다.

\subsection{무한 구간}

$[a,\infty)$의 적분은 $x=a+(1-t)/t$로 $(0,1]$의 적분으로 바꾼다. 야코비안 $1/t^2$이 곱해지므로 피적분함수가 무한대에서 충분히 빨리 줄어야 바뀐 함수가 $t=0$에서 유계이다. QUADPACK의 QAGI는 이 변환 뒤에 15점 Kronrod 규칙과 적응 분할, $\varepsilon$-알고리즘을 쓴다. $\int_{-\infty}^\infty e^{-x^2}dx$는 함수 호출 $270$번에 정확한 $\sqrt\pi$이다. 진동하며 천천히 주는 $\int_0^\infty\sin x/x\,dx=\pi/2$는 변환으로는 다루기 어렵다. QAWF는 홀수 개의 반주기로 이루어진 구간($\omega=1$이면 길이 $3\pi$)마다 QAWO로 적분해 부호가 번갈아 바뀌는 부분합을 만들고 그 부분합의 수열에 $\varepsilon$-알고리즘을 쓰는데, 함수 호출 $430$번에 오차 $1.8\times10^{-11}$을 낸다.

\section{이중 지수 구적}
\label{sec:tanh-sinh}

\subsection{변수 바꾸기}

35장에서 본 대로 사다리꼴 규칙은 실수축 전체에서 해석적이고 빨리 주는 함수에서 기하급수적으로 수렴한다. 그러므로 유한 구간의 적분을 그런 함수의 적분으로 바꾸면 된다. Takahasi와 Mori의 \term{이중 지수 변환}{double exponential transformation}은
\[
x=\tanh\Bigl(\frac\pi2\sinh t\Bigr),\qquad\frac{dx}{dt}=\frac{(\pi/2)\cosh t}{\cosh^2\bigl((\pi/2)\sinh t\bigr)}
\]
으로 $(-1,1)$을 실수축 전체로 보낸다\cite{takahasi1974}. 야코비안이 $t\to\pm\infty$에서 $\exp(-(\pi/2)e^{|t|})$처럼 이중 지수적으로 줄어서 끝점에서 $(1-x)^{-1/2}$이나 $\log(1-x)$처럼 발산하는 특이점이 있어도 바뀐 피적분함수는 여전히 이중 지수적으로 준다. 그 위에서 간격 $h$의 사다리꼴 규칙을 쓰면 오차가 대략 $\exp(-c/h)$로, 점의 개수 $N$에 대해 $\exp(-cN/\log N)$으로 준다\cite{mori2001}. 이것을 흔히 \term{tanh-sinh 구적}{tanh-sinh quadrature}이라 한다. 간격을 반으로 줄일 때 앞 단계의 점을 모두 다시 쓸 수 있으므로, 수준마다 점이 두 배가 되면서 이전 수준의 값을 버리지 않는다.

$[0,1]$의 네 적분에 이 변환을 쓰면, 수준 $0,1,2,3$(점 $13$, $25$, $49$, $97$개)에서 오차가 $x^{-1/2}$은 $4.5\times10^{-3}$, $4.6\times10^{-7}$, $3.1\times10^{-15}$, $2.2\times10^{-16}$, $\log x$는 $2.4\times10^{-2}$, $1.0\times10^{-5}$, $1.5\times10^{-13}$, $2.2\times10^{-16}$이고, 끝점에서 미분 불가능한 $\sqrt{1-x^2}$은 $3.5\times10^{-2}$, $2.9\times10^{-5}$, $1.5\times10^{-12}$, $0$이다. 수준마다 맞는 자릿수가 두 배가 되는 모양이다. 특이점이 없지만 극점이 가까운 $1/(1+25x^2)$은 수준 $3$에서 $1.2\times10^{-11}$로 오히려 조금 느리다. 이중 지수 변환은 극점을 $t$-평면의 실수축 가까이로 끌어오기 때문이다.

\subsection{끝점 근처의 반올림}

변환된 노드는 끝점에 이중 지수적으로 가까이 몰린다. 수준이 높으면 $1-x$가 $10^{-300}$보다 작은 노드가 생기는데, 이것을 $x$로 저장한 뒤 $1-x$를 계산하면 binary64에서 $0$이 되거나 소거로 정확도를 잃는다. 그래서 구현은 $x$ 대신 끝점까지의 거리 $1-x=1/\bigl(e^{u}\cosh u\bigr)$, $u=(\pi/2)\sinh t$를 직접 계산해 넘기고 피적분함수가 끝점까지의 거리를 받을 수 있게 하거나 $x$가 끝점과 같아진 노드는 무게를 $0$으로 둔다. Bailey 등은 높은 정밀도의 적분에서 이 방법이 Gauss 구적이나 오차함수 변환보다 대개 효율적임을 보였고\cite{bailey2005}, Vanherck 등은 부동소수점에서의 세부를 정리했다\cite{vanherck2020}. scipy의 \code{tanhsinh}가 그들을 따른다.

\section{다차원 적분}
\label{sec:multidim}

여러 변수의 적분을 한 변수 적분의 반복으로 쓰면 가장 안쪽부터 \code{quad}를 겹쳐 부를 수 있다. 차원이 $d$이고 한 변수에 $m$번씩 호출하면 $m^d$번이 든다. 오차의 어림도 문제이다. 안쪽 적분의 오차가 바깥 적분의 피적분함수에 잡음으로 들어가므로, 층마다의 어림 가운데 가장 큰 것을 내는 흔한 방식은 엄밀한 상계가 아니다.

영역 전체에 다차원 규칙을 쓰고 오차가 큰 영역을 나누는 방법도 있다. 한 변수의 Gauss--Kronrod를 축마다 곱한 규칙은 점이 $21^d$개여서 차원이 조금만 커도 감당할 수 없다. Genz와 Malik의 규칙은 $d$차원 상자에서 차수 $7$까지 정확하면서 점을 $2^d+2d^2+2d+1$개만 쓰고 차수 $5$의 내장 규칙으로 오차를 어림한다\cite{genz1980}. 차원이 열을 넘으면 이것도 $2^d$ 때문에 어렵고 60장의 몬테카를로와 준몬테카를로가 남는다.

\section{급수의 합}
\label{sec:series}

무한 급수 $\sum_{k\ge a}f(k)$에서 $f$가 매끄럽고 단조 감소하면 적분 판정법의 부등식
\begin{equation}
\label{eq:integral-test}
\int_{n+1}^\infty f(x)\,dx\le\sum_{k>n}f(k)\le\int_n^\infty f(x)\,dx
\end{equation}
가 꼬리를 감싼다. $f$가 감소하므로 $[k,k+1]$에서 $f(k+1)\le f(x)\le f(k)$이다. 이것을 적분해 $k\ge n$에 대해 더하면 오른쪽 부등식이, $k>n$에 대해 더하면 왼쪽 부등식이 나온다. 꼬리의 적분은 앞의 tanh-sinh로 계산할 수 있으므로, 처음 $n$항을 직접 더하고 꼬리를 적분으로 어림하면 된다. Euler--Maclaurin 공식 \eqref{eq:euler-maclaurin}을 거꾸로 쓰면 더 정확하다. $\sum_{k>n}f(k)=\int_n^\infty f-f(n)/2-f'(n)/12+\cdots$이다. $\sum1/k^2=\pi^2/6$의 $1000$항까지의 부분합은 오차가 $1.0\times10^{-3}$이지만, 꼬리의 전개 $1/n-1/(2n^2)+1/(6n^3)$을 더하면 오차가 $1.6\times10^{-15}$이다. 교대급수에는 Cohen, Rodriguez Villegas, Zagier의 가속이 오차를 항마다 약 $5.8$배씩 줄인다\cite{cohen2000}.

\section{stardust의 적분}
\label{sec:quadrature2-stardust}

\begin{instardust}[\code{docs/design/integrate.md} §3.5]
stardust의 \code{tanhsinh}는 scipy의 구현을 원소마다 옮겼다. scipy는 노드와 무게의 표를 numpy의 \code{cosh}, \code{sinh}, \code{exp}로 계산하는데, 이 점들에서 그 값은 macOS의 수학 라이브러리와 같다. stardust도 C++ 표준 라이브러리의 같은 함수로 표를 만들어 수준 $0$에서 $10$까지의 $16386$개 값이 scipy와 비트까지 같다. 자체 배열 영역의 벡터화된 \code{exp}는 이 점들의 약 3분의 1에서 끝자리가 달라 쓰지 않았다. 설계 문서에 따르면 이 표 자체는 정확하지 않다. 같은 식의 참값과 견주면 최대 $980$ulp, 평균 $47$ulp 틀리는데, $u_2$의 반올림 오차가 $e^{u_2}$에서 $u_2$배로 커지기 때문이다. 꼬리의 무게가 작아 적분값에는 거의 영향이 없으므로, scipy와 같은 판정을 내리는 쪽을 택해 표를 그대로 쓴다. 수준마다의 합도 멈춤 판정을 정하므로 numpy의 쌍별 합 순서를 그대로 따른다.
\end{instardust}

\begin{instardust}[\code{docs/design/integrate.md} §3.2, §3.6]
\code{nquad}는 가장 안쪽부터 \code{quad}를 겹치고 돌려주는 오차는 scipy처럼 모든 층의 어림 가운데 최댓값이다. 설계 문서는 처음에 이것을 바깥 층의 어림으로 잘못 읽었다가 바로잡았고 그것이 엄밀한 상계가 아니라는 점을 API 문서에 적었다. 급수의 합 \code{nsum}은 \eqref{eq:integral-test}의 생각대로, 꼬리의 적분을 \code{tanhsinh}로 구하고 항이 허용오차보다 작아지는 곳까지 직접 더한 뒤 그 경계에서 적분으로 넘어간다. $\sum1/k^2$에서 함수를 $10$번 불러 $8561$개의 항을 계산하고 실제 오차 $3\times10^{-13}$에 어림 $7.4\times10^{-9}$을 낸다.
\end{instardust}

\section{맺음}

적분 루틴의 실력은 규칙 하나의 정확도보다 함수의 나쁜 부분을 다루는 방법에서 드러난다. 적응 분할은 나쁜 곳에 노드를 모으고, 수열 가속은 나쁜 곳을 나누며 생긴 규칙적인 오차를 외삽해 지우며, 무게 함수와 변수 바꾸기는 나쁜 부분을 아예 해석적으로 처리하거나 보이지 않게 만든다. 어느 방법이든 함수를 상자로 보는 이상 놓치는 봉우리는 놓친다는 한계는 남는다. 다음 장에서는 이 장과 앞 장에서 여러 번 등장한 주기 함수와 사다리꼴 규칙의 이야기를 푸리에 해석으로 정리한다.

적응 적분과 QUADPACK의 설계는 Piessens 등의 책\cite{piessens1983quadpack}에, 이중 지수 변환의 이론은 Mori와 Sugihara의 개관\cite{mori2001}에 정리되어 있다.

\chapterreferences
